\section{Numerical methods and sensitivity analyses}\label{si:numerics}

\paragraph{Numerical method.} The delay equations are integrated by the method of steps with an explicit first-order Euler scheme. At step $k$, the current states are read at $t_{k-1}$; time-dependent and delayed inputs are evaluated at the next grid time $t_k$, with linear interpolation of delayed histories. The quality velocity $G_c(q(t_k-\tau_g))$ is projected onto $T_{K_q}(q(t_{k-1}))$ before the Euler update; the updated quality is then clipped to $K_q=[0,q_{\max}]$ as a finite-step safeguard. The remaining states are advanced using the computed rates, their stated bounds are applied, reserve area is calculated as $A_r=A_{tot}-A_f-A_c$, and the updated $B$ is stored in the delay history. The positive-part operator is exact. For ordinary runs, pre-run histories $q(s)=q_0$ and $B(s)=B_0$ are constant over $[-\max(\tau_g,\tau_p),0)$. In the technology-wave experiment, both runs use the common no-wave prehistory; the active $T_b(0)$ changes current $B(0)$, and subsequent wave-adjusted $B$ values enter the delayed history used for $K(t-\tau_p)$. The demonstration run and bisection thresholds use $\Delta t=0.05$\,yr; drawdown-depth sweeps use $\Delta t=0.1$\,yr, with selected thresholds evaluated at $0.05$\,yr. For $A_{c0}=1.0$\,ha, the $10^{-6}$\,ha floor-proximity threshold changes from $2.955$ to $2.98$\,cap across these step sizes; threshold values are therefore reported as tolerance- and resolution-dependent bands.

\paragraph{Delayed versus current quality in the regeneration term.} The reported model uses the projected delayed law $\dot q=\Pi_{T_{K_q}(q(t))}G_c(q(t-\tau_g))$. A current-quality comparator, using $G_c(q(t))$ with the same imposed interval, is retained only as a sensitivity check. The quality index is uncalibrated, $q_{\max}=1$ is an imposed admissibility bound, and $\rho_c$ and $\tau_g$ are scenario settings; the comparison is not a validation test. In the baseline history $q(s)=0.9$ with $\rho_c=0.08$\,yr$^{-1}$, the projected paths for $\tau_g=20$ and $40$\,yr both reach $q_{\max}$ at $125/9\approx13.8889$\,yr, before either delay elapses. Their equality is a saturation-induced degeneracy, not general delay robustness. The current-quality comparator extends the demonstration window by $1.3$\,yr; the new floor-threshold run changes by at most $0.001$\,cap across the displayed finite set of delays (Table~\ref{tab:si_delay}).

\begin{table}[htbp]
\centering\small
\caption{Sensitivity to evaluating regeneration at current quality rather than using the projected delayed law. The projected delayed model is the specified model; current quality is a comparator only. Demonstration summaries use $\Delta t=0.05$\,yr and $T=400$\,yr; floor-proximity thresholds use $\Delta t=0.1$\,yr and $T=800$\,yr, with $\tau_p=25$\,yr. Biocapacity and contributions are in gha, land areas in ha, $P_0$ in cap, and delays and step intervals in yr.}
\label{tab:si_delay}
\begin{tabularx}{\textwidth}{>{\raggedright\arraybackslash}Xcc}
\toprule
Quantity & current $q(t)$ & projected delayed $q(t-\tau_g)$ \\
\midrule
Demonstration window with $\dot B>0$, $\dot A_c<0$ (yr) & $150.3$ & $149.0$ \\
Mask-end values $B$, $A_c$, $A_f$ & $0.825,\ 0.808,\ 0.923$ & $0.825,\ 0.808,\ 0.923$ \\
Capital contribution $\Ups(q)A_c$ (initial $\to$ mask end, gha) & $0.08364\to0.04038$ & $0.08364\to0.04038$ \\
Fast-land contribution $Y_f$ (initial $\to$ mask end, gha) & $0.51000\to0.78462$ & $0.51000\to0.78462$ \\
$P_0$ floor threshold ($A_{c0}=1.0$\,ha, $\tau_g=5$\,yr, $\Delta t=0.1$\,yr) & $2.955$ & $2.955$ \\
$P_0$ floor threshold ($A_{c0}=1.0$\,ha, $\tau_g=20$\,yr, $\Delta t=0.1$\,yr) & $2.955$ & $2.954$ \\
$P_0$ floor threshold ($A_{c0}=1.0$\,ha, $\tau_g=40$\,yr, $\Delta t=0.1$\,yr) & $2.955$ & $2.954$ \\
\bottomrule
\end{tabularx}

\end{table}

For $\tau_g=20$ and $40$\,yr, the constant history drives $q$ to the cap before the delayed state first enters the solution, so the quality paths are identical; matching thresholds in those two rows are not evidence that delay is generally unimportant. The values in the table were regenerated for this revision from the v45-derived projected solver (thresholds: $T=800$\,yr, $\Delta t=0.1$\,yr). Their agreement at displayed precision with archived C10 values does not identify the historical code-to-output path.

\paragraph{Well-posedness.} For $\rho_c\ge0$, $q_{\max}>0$, and any admissible continuous history $\phi:[-\tau_g,0]\to K_q$, the quality rate $G_c(\phi)$ is nonnegative and Lipschitz on $K_q$. Before the first hit of $q_{\max}$, the method-of-steps delayed equation therefore has a unique solution; at the upper boundary, tangent-cone projection sets the outward velocity to zero, and the state remains there. The lower bound is invariant because the rate is nonnegative. The remaining model right-hand sides are locally Lipschitz on each stated active set; the positive-part and minimum operators are Lipschitz, and the land/debt/population bounds are enforced as described above. We make no smooth linear-stability claim across the quality ceiling or the shortfall switch; the affine spectrum below is a one-sided deficit-branch result.

\section{Extended model results}\label{si:results}

\subsection{Fast-land yield sweep and parameter sensitivity of the composition diagnostic}\label{si:sweep}

Table~\ref{tab:si_sweep} extends the fast-land yield sweep in Table~2(a) in the main text with the slope of the balance branch. The sign of $\dconv$ remains positive as $b_{G,c}$ varies threefold: at $b_f=0.85$\,gha\,ha$^{-1}$, $\Ups(q^*)=0.058,\,0.066,\,0.074$\,gha\,ha$^{-1}$ for $b_{G,c}=0.4,\,0.8,\,1.2$\,gha\,yr\,ha$^{-1}$, so $\dconv=+0.792,\,+0.784,\,+0.776$\,gha\,ha$^{-1}$.

\begin{table}[htbp]
\centering\small
\caption{Fast-land yield sweep at the baseline (all other parameters fixed): composition diagnostic at $q^*$, eigenvalues of the deficit-side matrix, slope of the balance branch and demand ceiling.}
\label{tab:si_sweep}
\begin{tabular}{l r r r r r}
\toprule
$b_f$ (gha\,ha$^{-1}$) & $\dconv(q^*)$ (gha\,ha$^{-1}$) & $\lambda_1$ (yr$^{-1}$) & $\lambda_2$ (yr$^{-1}$) & $b_fk-b_c$ (gha\,ha$^{-1}$) & $\Eceil$ (gha) \\
\midrule
$0.60$ & $+0.534$ & $-0.0522$ & $-0.975$ & $0.277$ & $0.804$ \\
$0.85$ & $+0.784$ & $-0.0537$ & $-0.998$ & $0.414$ & $1.138$ \\
$1.20$ & $+1.134$ & $-0.0547$ & $-1.014$ & $0.605$ & $1.606$ \\
$2.00$ & $+1.934$ & $-0.0556$ & $-1.029$ & $1.041$ & $2.675$ \\
\bottomrule
\end{tabular}
\end{table}

\subsection{Drawdown-depth ladder}\label{si:ladder}

Table~\ref{tab:si_drawdown} gives the minimum capital land reached during endogenous-demand runs ($A_{f0}=0.6$, $q_0=0.9$, $\tau_g=20$, $\tau_p=25$\,yr, $\chi_r=0.10$, $\Delta t=0.1$\,yr, $T=400$\,yr) on a grid of initial population $P_0$; the full grid (step $0.1$) is in \texttt{si\_drawdown\_depth.csv}. The minimum is reached when the shortfall first closes ($B\ge E$), between $3$ and $25$\,yr after the start depending on $P_0$. At $A_{c0}=1.00$\,ha, the $\tau_g=20$ and $40$\,yr runs have identical quality paths because the baseline history reaches the imposed ceiling before either delay elapses; their matching threshold is saturation-induced degeneracy. The $\tau_g=5$\,yr case is distinct. The finite delay comparison does not establish general delay robustness (SI Table~\ref{tab:si_delay} in \S~\ref{si:numerics}). The depth is continuous and monotone in $P_0$ and in $A_{c0}$. The initial population $P_0$ (cap) at which the minimum comes within $10^{-2}$, $10^{-3}$ and $10^{-4}$\,ha of the floor is $1.8,\,1.9,\,2.1$ ($A_{c0}=0.40$\,ha); $2.2,\,2.3,\,2.4$ ($A_{c0}=0.70$\,ha); $2.6,\,2.7,\,2.8$ ($A_{c0}=1.00$\,ha); $2.8,\,2.9,\,3.0$ ($A_{c0}=1.18$\,ha); the bisection with a $10^{-6}$\,ha floor-proximity tolerance ($\Delta t=0.05$\,yr, $T=1000$\,yr) gives $2.41,\,2.70,\,2.98,\,3.14$\,cap. \par\pagebreak

The thresholds are identical to three decimals for $\chi_r=0.10$ and $\chi_r=0$ ($2.387$ and $2.955$\,cap at $A_{c0}=0.40$ and $1.00$\,ha, $\Delta t=0.1$\,yr).

\begin{table}[htbp]
\centering\small
\caption{Minimum capital land $\min_tA_c$ (ha) reached during the run, by initial capital $A_{c0}$ (rows, ha) and initial population $P_0$ (columns, cap); typed floor $\Acmin=0.05$\,ha.}
\label{tab:si_drawdown}
\begin{tabular}{lrrrrrrrr}
\toprule
$A_{c0}$ $\backslash$ $P_0$ & 1.0 & 1.5 & 2.0 & 2.5 & 2.8 & 3.0 & 3.5 & 4.0 \\
\midrule
$0.4$ & 0.396 & 0.147 & 0.050 & 0.050 & 0.050 & 0.050 & 0.050 & 0.050 \\
$0.7$ & 0.700 & 0.442 & 0.121 & 0.050 & 0.050 & 0.050 & 0.050 & 0.050 \\
$1.0$ & 1.000 & 0.731 & 0.398 & 0.075 & 0.050 & 0.050 & 0.050 & 0.050 \\
$1.18$ & 1.180 & 0.898 & 0.557 & 0.230 & 0.055 & 0.050 & 0.050 & 0.050 \\
\bottomrule
\end{tabular}

\end{table}

\subsection{Recovery metric and the role of the restoration flow}\label{si:recovery}

The recovery metric is $\mathcal R_c(T)=(A_c(T)-A_c^{\mathrm{deg}})/(A_c^{\mathrm{init}}-A_c^{\mathrm{deg}})$, with $\mathcal R_c=0$ no recovery, $1$ full recovery to the reference, $<0$ continued decline and $>1$ over-restoration. Table~\ref{tab:si_recovery} evaluates it on the projected delayed-quality drawdown runs with $A_{c0}=1.0$\,ha ($A_c^{\mathrm{init}}=1.0$\,ha, $A_c^{\mathrm{deg}}=\min_tA_c$ in ha, $T=400$\,yr). Without restoration the capital-land area remains at its minimum because no return flow exists. With the surplus-gated flow, recovery is partial across these six initial-population cases ($\mathcal R_c=0.21$--$0.89$); none reaches the reference. The distinct gate-sign experiment in SI~\S S4 tests other initial states and intervention settings, including cases of over-restoration.

\begin{table}[htbp]
\centering\small
\caption{Recovery of the capital-land area after a demand-driven drawdown ($A_{c0}=1.0$\,ha, $A_{f0}=0.6$\,ha, $q_0=0.9$, $T=400$\,yr): initial population $P_0$ is in cap, initial footprint $E_0$ in gha, capital-land values in ha, and the recovery metric is dimensionless.}
\label{tab:si_recovery}
\begin{tabular}{l r r r r r r}
\toprule
$P_0$ & $E_0$ & $\min_tA_c$ & $t(\min)$ (yr) & $A_c(T)$, $\chi_r=0$ & $A_c(T)$, $\chi_r=0.10$ & $\mathcal R_c(T)$, $\chi_r=0.10$ \\
\midrule
$1.5$ & $0.825$ & $0.731$ & $3.0$ & $0.731$ & $0.969$ & $0.89$ \\
$2.0$ & $1.100$ & $0.398$ & $3.8$ & $0.398$ & $0.677$ & $0.46$ \\
$2.5$ & $1.375$ & $0.075$ & $4.1$ & $0.075$ & $0.433$ & $0.39$ \\
$3.0$ & $1.650$ & $0.050$ & $13.8$ & $0.050$ & $0.328$ & $0.29$ \\
$3.5$ & $1.925$ & $0.050$ & $20.6$ & $0.050$ & $0.278$ & $0.24$ \\
$4.0$ & $2.200$ & $0.050$ & $26.6$ & $0.050$ & $0.250$ & $0.21$ \\
\bottomrule
\end{tabular}

\end{table}

\subsection{Gate-sign phase dependence}\label{si:gates}

The gate-sign comparison begins from a post-drawdown state ($A_c=0.058$\,ha near the floor, $A_f=0.6$\,ha, $q=0.9$, and $P=1.2$\,cap). At a degraded equilibrium with $R_B=1$, $B=E$, so neither the surplus-gated nor deficit-gated restoration flow is active. The reference equilibrium has $A_c^{\mathrm{init}}=0.8225$\,ha and $P^{\mathrm{init}}=1.489$\,cap; the degraded value is $A_c^{\mathrm{deg}}=0.0513$\,ha. All rows use $T_b=0$ and the stated initial state.

\begin{table}[H]
\centering\scriptsize
\setlength{\tabcolsep}{2pt}
\caption{Restoration gate and recovery outcomes at $T=100$ and $900$\,yr. The table reports end-state land areas, population relative to the reference, trajectory extrema, and peak debt $D$ (gha\,yr). Annual account values $B=E$ are in gha; land areas are in ha; population ratios and recovery metrics are dimensionless. Each row approaches $R_B=1$ ($E=B$).}
\label{tab:si_gates}
\begin{tabular}{l r r r r r r r r r r}
\toprule
Policy / gate & $\mathcal R_c(100)$ & $\mathcal R_c(900)$ & $A_c$ & $A_f$ & $A_r$ & $B=E$ & $P/P^{\mathrm{init}}$ & $\min A_r$ & $\max A_c$ & $D_{peak}$ \\
\midrule
none ($R_{rc}=0$) & $-0.00$ & $-0.00$ & $0.051$ & $1.336$ & $1.113$ & $1.138$ & $1.39$ & $1.113$ & $0.057$ & $0.152$ \\
surplus-gated, $\chi_r=0.02$ & $0.69$ & $0.73$ & $0.612$ & $1.030$ & $0.858$ & $0.906$ & $1.11$ & $0.858$ & $0.612$ & $0.152$ \\
surplus-gated, $\chi_r=0.05$ & $1.08$ & $1.07$ & $0.875$ & $0.886$ & $0.738$ & $0.797$ & $0.97$ & $0.724$ & $0.885$ & $0.152$ \\
surplus-gated, $\chi_r=0.10$ & $1.28$ & $1.27$ & $1.030$ & $0.802$ & $0.668$ & $0.733$ & $0.90$ & $0.602$ & $1.092$ & $0.152$ \\
surplus, $\chi_r=0.10$, target cap $A_c^*=0.8225$ & $1.00$ & $1.00$ & $0.822$ & $0.915$ & $0.763$ & $0.819$ & $1.00$ & $0.664$ & $0.822$ & $0.152$ \\
surplus, $\chi_r=0.10$, conversion frozen & $1.35$ & $1.36$ & $1.097$ & $0.765$ & $0.638$ & $0.705$ & $0.86$ & $0.602$ & $1.097$ & $0.543$ \\
deficit-gated, $\chi_r=0.05$ & $0.00$ & $0.00$ & $0.055$ & $1.334$ & $1.112$ & $1.136$ & $1.39$ & $1.112$ & $0.060$ & $0.145$ \\
deficit-gated, $\chi_r=0.10$ & $0.01$ & $0.01$ & $0.059$ & $1.331$ & $1.110$ & $1.135$ & $1.39$ & $1.110$ & $0.065$ & $0.139$ \\
deficit, $\chi_r=0.10$, conversion frozen & $0.05$ & $0.05$ & $0.088$ & $1.315$ & $1.096$ & $1.123$ & $1.37$ & $1.096$ & $0.088$ & $0.159$ \\
\bottomrule
\end{tabular}
\end{table}

Without restoration, $A_c$ remains at its minimum because the model has no return flow. Surplus-gated restoration yields partial or over-restoration depending on $\chi_r$; at $\chi_r\ge0.05$, the flow continues while $B>E$, and the supported population falls to $0.97$--$0.90$ of its reference value. A target cap returns $A_c$ to its reference value ($\mathcal R_c=1.00$). Deficit-gated restoration acts during the shortfall, but at the tested rates its contribution is smaller than concurrent conversion and the recovery fraction remains near zero. Freezing conversion under surplus-gated restoration raises peak debt from $0.152$ to $0.543$\,gha\,yr; restoration then reallocates reserve area to capital land while fast-land maintenance is reduced. These outcomes depend on the specified initial state and intervention rates. The gate-sign comparisons use $T_b=0$; the controlled wave/no-wave debt comparison is reported in SI~\S S8.\par\pagebreak

\subsection{Regime table}\label{si:regimes}

\begin{table}[H]
\centering\small
\caption{The composition diagnostic $\dconv$ (gha\,ha$^{-1}$) as the decision condition for mask visibility. There is no interior area-based capacity maximum in the quality-based model; the only shape distinction is the sign of $\dconv$.}
\label{tab:si_regime}
\begin{tabularx}{\textwidth}{l l X}
\toprule
Regime & Condition & Behavior \\
\midrule
composition-premium & $\dconv>0$ & the conversion contribution to accounting $B$ is positive; net $B$ also depends on other flows (the mask is possible in the baseline) \\
composition-neutrality & $\dconv=0$ & first-order neutrality; non-generic \\
composition-discount & $\dconv<0$ & conversion lowers $B$ at the conversion margin; conversion alone cannot produce a composition mask \\
\bottomrule
\end{tabularx}
\end{table}

\section{Empirical details}\label{si:empirical}

\subsection{Data and vintages}\label{si:vintages}

The empirical backbone is the National Ecological Footprint and Biocapacity Accounts (NFA), Edition 2026, results and metadata version 1.0 \citep{Lo2026NFA}. Although the workbook covers 1961--2025, the primary analysis uses complete world observations for 1961--2023; the 2024--2025 estimated rows are excluded. The 1961--2022 sensitivity uses the same edition and population series. The release recalculates common historical years and reports biocapacity by fishing grounds, built-up land, cropland, grazing land, and forest biocapacity. These are account categories, not a common-boundary physical-stock inventory.

The workbook's $\texttt{pop\_1000\_persons}$ field is UN PopStat population in thousands of persons; it is converted to persons before recomputing per-capita values. The field $\texttt{efp\_total\_gha}$ is the Global Ecological Footprint of Production, denoted $E_{\mathrm{prod}}$ here. Per-capita values are ratios of the unrounded account totals and population, not the workbook's rounded per-capita columns. Table~\ref{tab:si_nfa_endpoints} summarizes endpoint changes. For 1961--2023, $B$ is $9{,}725{,}417{,}562$ to $12{,}142{,}843{,}961$\,gha, population is $3.06487$ to $8.09173$ billion, $B/P$ is $3.17319$ to $1.50065$\,gha\,person$^{-1}$, and $E_{\mathrm{prod}}/P$ is $2.19700$ to $2.56462$\,gha\,person$^{-1}$. The identity $\Delta\ln B=\Delta\ln P+\Delta\ln(B/P)$ holds to rounding; it is an extensive/intensive accounting split, not a causal decomposition. The first year in which $E_{\mathrm{prod}}/B>1$ is 1973 in both windows.

\begin{table}[htbp]
\centering\small
\caption{World endpoint changes in NFA Edition 2026. The 1961--2022 column is a same-edition sensitivity; 2024--2025 estimates are excluded. Population and per-capita quantities are recomputed from workbook totals.}
\label{tab:si_nfa_endpoints}
\begin{tabularx}{\textwidth}{>{\raggedright\arraybackslash}Xrr}
\toprule
Measure & 1961--2023 & 1961--2022 \\
\midrule
$\Delta B/B_{1961}$ & $+24.8568\%$ & $+24.5450\%$ \\
$\Delta\ln B$ & $+0.221997$ & $+0.219497$ \\
$\Delta P/P_{1961}$ & $+164.0154\%$ & $+161.7211\%$ \\
$\Delta\ln P$ & $+0.970837$ & $+0.962109$ \\
$\Delta(B/P)/(B/P)_{1961}$ & $-52.7085\%$ & $-52.4131\%$ \\
$\Delta\ln(B/P)$ & $-0.748840$ & $-0.742612$ \\
$\Delta(E_{\mathrm{prod}}/P)/(E_{\mathrm{prod}}/P)_{1961}$ & $+16.7327\%$ & $+17.6994\%$ \\
First year $E_{\mathrm{prod}}/B>1$ & 1973 & 1973 \\
\bottomrule
\end{tabularx}
\end{table}

\newpage
\subsection{Divisia decomposition, sensitivity to the base-year share and level fallback}\label{si:divisia}

The crop proxy is an illustrative accounting construction. Let $A_f$ be the independently sourced OWID/FAOSTAT world cropland-area series (arable plus permanent-crop area, ha) and let $\tilde b_f$ be the independently sourced cereal-yield series (tonnes per hectare) \citep{OWIDCereal2026,OWIDCropland2026}. Cereal yield is a sentinel for crop productivity, not a measure of all-crop output. For a declared 1961 account-share anchor $\phi_0$, define $\hat b_f(t)=\hat b_f(1961)\tilde b_f(t)/\tilde b_f(1961)$, where $\hat b_f(1961)A_f(1961)=\phi_0B(1961)$, then $X_t=\hat b_f(t)A_f(t)$ and $C_t=B_t-X_t$. Thus $X$ is an index-scaled account proxy and $C$ is the residual required to recover the NFA aggregate; $C$ is not observed ecological capital, non-cropland biocapacity, or the model quantity $\Ups(q)A_c$.

When $X$ and $C$ are positive at both endpoints, the exact log-mean Divisia identity is
\begin{equation}
\Delta\ln B=w_X\bigl(\Delta\ln\tilde b_f+\Delta\ln A_f\bigr)+w_C\Delta\ln C,
\qquad w_X=\frac{L(X_T,X_0)}{L(B_T,B_0)},\quad w_C=\frac{L(C_T,C_0)}{L(B_T,B_0)},
\end{equation}
where $L(a,b)=(b-a)/(\ln b-\ln a)$ is the logarithmic mean, with its continuous limit at $a=b$. The measured NFA cropland share in 1961 is $\phi_0=0.1286472202$ and is used as the central anchor, but the mapping of that account share to the external area--yield proxy is an assumption, not an independently identified physical-capital observation. Table~\ref{tab:si_phi0} reports the full selected sensitivity for 1961--2023. At the central anchor, the cereal-yield, area, and residual terms are $+0.268662$, $+0.037710$, and $-0.084375$, summing to $\Delta\ln B=+0.221997$. The same-edition 1961--2022 sensitivity gives $+0.264696$, $+0.037506$, and $-0.082705$, summing to $+0.219497$. The decomposition is exact for the defined proxy but its allocation depends on $\phi_0$ and must not be read as a physical-stock attribution.

\begin{table}[htbp]
\centering\footnotesize
\caption{Sensitivity of the proxy Divisia decomposition to the 1961 account-share anchor $\phi_0$ (world, 1961--2023). Contributions and weights are dimensionless. The $\phi_0=0.128647$ row uses the measured NFA cropland share in 1961; other rows are illustrative calibrations.}
\label{tab:si_phi0}
\begin{tabular}{r r r r r r r}
\toprule
$\phi_0$ & $w_X$ & $w_C$ & Cereal yield & Area & Residual $C$ & Sum \\
\midrule
$0.100000$ & $0.1833$ & $0.7957$ & $+0.2088$ & $+0.0293$ & $-0.0162$ & $+0.2220$ \\
$0.128647$ & $0.2358$ & $0.7352$ & $+0.2687$ & $+0.0377$ & $-0.0844$ & $+0.2220$ \\
$0.150000$ & $0.2749$ & $0.6893$ & $+0.3133$ & $+0.0440$ & $-0.1352$ & $+0.2220$ \\
$0.190000$ & $0.3483$ & $0.6008$ & $+0.3968$ & $+0.0557$ & $-0.2305$ & $+0.2220$ \\
$0.250000$ & $0.4582$ & $0.4581$ & $+0.5221$ & $+0.0733$ & $-0.3734$ & $+0.2220$ \\
$0.300000$ & $0.5499$ & $0.3178$ & $+0.6265$ & $+0.0879$ & $-0.4924$ & $+0.2220$ \\
$0.320000$ & $0.5866$ & $0.2454$ & $+0.6683$ & $+0.0938$ & $-0.5401$ & $+0.2220$ \\
$0.340000$ & $0.6232$ & $0.1009$ & $+0.7100$ & $+0.0997$ & $-0.5877$ & $+0.2220$ \\
\bottomrule
\end{tabular}
\end{table}

For $\phi_0\ge 0.340531$ in 2023 (threshold $0.343400$ for 2022), the residual $C$ is nonpositive at the terminal endpoint and logarithmic weighting is invalid. In that range the exact level identity $\Delta B=\Delta X+\Delta C$ is used instead; it remains an accounting identity and carries no stock interpretation. The selected anchor has $X_{2023}/B_{2023}=0.377784$ and a positive residual, but the proxy fast-land share and residual contribution vary across the sensitivity values. These calculations use the same NFA $B$ in all rows; the residual absorbs the difference between the index-scaled crop proxy and that aggregate, including account-category, boundary, valuation, vintage, and calibration differences.

\subsection{Cropland-account calibration}\label{si:cropland}

Direct NFA cropland-account movements are reported separately from the illustrative proxy. From 1961 to 2023, NFA cropland biocapacity rises $2.870055$-fold while NFA cropland area rises $1.198029$-fold; the quotient of account biocapacity by NFA area rises $2.395647$-fold. For the same-edition 1961--2022 sensitivity the ratios are $2.851441$, $1.198326$, and $2.379520$. The per-hectare quantity is an NFA account-value quotient, not an independent physical crop-yield measure. By comparison, the external OWID/FAOSTAT area series rises $1.173414$-fold and cereal yield rises $3.124667$-fold over 1961--2023. Those independent series provide a sentinel proxy, not a reconstruction of the NFA cropland account or a measure of all-crops output. The model's representative $b_{f0}=0.85\,\mathrm{gha\,ha^{-1}}$ remains a conditional parameter, not an empirical fit to these account ratios.

\subsection{Observation-operator hierarchy}\label{si:operator}

\begin{itemize}[leftmargin=1.5em]
\item \emph{NFA accounting observations.} Total and component biocapacity values in gha support statements about the composition and movement of the NFA accounts under the stated edition and conventions. They do not directly observe biomass, soil condition, wetland extent, fish biomass, or the model state $A_c$.
\item \emph{Fast-provisioning area.} The external world cropland-area series is arable plus permanent-crop area, 1961--2023 (ha) \citep{OWIDCropland2026}. It is used as an independently sourced proxy input, not treated as identical to the NFA cropland-area series.
\item \emph{Crop-productivity sentinel.} The OWID/FAOSTAT cereal-yield series, 1961--2024, is measured in tonnes per hectare \citep{OWIDCereal2026}. It is not an all-crops index and is not converted into a physical stock. The base-anchored $\hat b_f$ is only the index-scaled term in the illustrative account proxy.
\item \emph{Ecological-capital components.} Forest extent, growing stock or biomass, soil carbon and condition, wetland extent and condition, and species- or stock-specific fish biomass could inform separate physical observation operators. They would require harmonized spatial boundaries, reference years, definitions, uncertainty, and explicit treatment of overlap. No common-boundary global physical-stock series for the relevant components was assembled here. Existing NFA forest biocapacity and the residual $C$ are not substitutes for those measurements.
\item \emph{Model mapping.} In the model, $Y_c=\Ups(q)A_c$ is a model-defined capital contribution. Empirically, neither $A_c$ nor $q$ is identified by $B$ or by an NFA component; the proxy residual $C=B-X$ is not assumed to equal $Y_c$. A defensible observation operator would need independent physical measures and a declared mapping before a model-to-data comparison.
\end{itemize}

Table~\ref{tab:si_model_nfa} summarizes this mapping without equating an NFA category to a model state.

\begin{table}[htbp]
\centering\small
\caption{Conceptual mapping from two-land model quantities to NFA information and independent proxy inputs. The rows specify possible observation operators, not identities between NFA categories and model states. Annual account values are in gha; land areas are in ha.}
\label{tab:si_model_nfa}
\begin{tabularx}{\textwidth}{>{\raggedright\arraybackslash}p{0.23\textwidth} >{\raggedright\arraybackslash}p{0.31\textwidth} X}
\toprule
Model quantity & Closest account or proxy information & Interpretation and identification limit \\
\midrule
Fast area and proxy account $A_f$, $X=\hat b_fA_f$ & NFA cropland account; independently sourced OWID cropland area and cereal yield & NFA cropland biocapacity is an account value; the external area--yield product is an illustrative, base-anchored proxy. Neither is an observed model fast-land yield or a one-to-one physical measure. \\
Capital area, quality and contribution $A_c$, $q$, $Y_c=\Ups(q)A_c$ & No direct NFA counterpart; the proxy residual $C=B-X$ is defined by subtraction & In the model $Y_c=\Ups(q)A_c$; empirically $C$ is not assumed equal to $Y_c$. It also absorbs differences in boundaries, account categories (including built-up land), valuation, vintage and proxy calibration, and is not a physical-stock measure. \\
Reserve area $A_r$ & No direct NFA account counterpart & $A_r$ closes the model's conserved-area ledger; it is not recoverable from the NFA aggregate. \\
\bottomrule
\end{tabularx}
\end{table}

\begin{table}[htbp]
\centering\small
\caption{Observability ladder: what is directly measured, what is identifiable by construction, and what remains unidentified.}
\label{tab:si_ladder}
\begin{tabularx}{\textwidth}{>{\raggedright\arraybackslash}p{0.3\textwidth} X >{\raggedright\arraybackslash}p{0.3\textwidth}}
\toprule
Tier & Quantities & Status \\
\midrule
Observable physical area input & OWID cropland area [ha]; other mapped land-cover classes where available & Measured extent, but mapping to effective $A_c$ is not direct \\
Observable population & $P$ [persons] & Direct in the UN PopStat workbook field after unit conversion \\
Observable productivity sentinel & Cereal yield [tonnes ha$^{-1}$] & Physical yield of cereals only; not all-crops output or ecological capital \\
Observable annual account value & NFA $B$, $E_{\mathrm{prod}}$ and components [gha] & Measured only under the NFA accounting convention; value-weighted accounts \\
Proxy account and residual & $X=\hat b_fA_f$ and $C=B-X$ [gha] & Identified by construction for the stated inputs and anchor; not physical stocks and not assumed equal to $\Ups(q)A_c$ \\
Effective capital contribution & $Y_c=\Ups(q)A_c$ [gha] & A model quantity; no common-boundary physical observation operator is assembled here \\
Not separately identifiable & $A_c$, $q$, $b_c$, $b_{G,c}$, $\rho_c$ and $D$ & Requires independent physical-state, quality, regeneration or degradation observations and a declared mapping \\
\bottomrule
\end{tabularx}
\end{table}

\subsection{Regeneration anchors}\label{si:anchors}

The capital-land regeneration timescale is a lumped recruitment-response timescale. Plot-level age--above-ground-biomass records from Neotropical secondary forests (2ndFOR; 1,334 plots, 41 chronosequences, 35 good fits) yield median $t_{50}\approx30$\,yr and $\tau_r\approx38$\,yr under a lagged saturating recovery model \citep{Poorter2016}. Deplete-then-rebuild trajectories for 33 collapsed fish stocks in RAM Legacy ($B/B_{MSY}$) yield $t_{50}\approx5.6$\,yr and $\tau_r\approx3.4$\,yr \citep{HutchingsReynolds2004,Neubauer2013}; soil organic carbon re-equilibrates over approximately 17--23\,yr following land-use change \citep{Poeplau2011}. Forest biomass reaches about half of old-growth values by 20\,yr and 90\% at a median of about 66\,yr, while approximately 29\% of collapsed fish stocks recover to 50\% within 5--15\,yr. These records span distinct systems and are not pooled into a single empirical rate. Forest and soil observations provide decadal-scale context for the representative effective response scale of about 29\,yr; the 10--40\,yr interval is a modeling range for the lumped capital-land class, not a pooled empirical estimate. The faster fish-stock episodes are not treated as a direct lower bound on that class. Neither $\rho_c$ nor the delay $\tau_g$ is calibrated from these observations.

\subsection{Recursive identifiability in the land-type accounts}\label{si:recursive}

NFA Edition 2026 shows why the aggregate should be read together with its account composition, without translating account movement into physical-stock movement. From 1961 to 2023, cropland contributes $+24.0577$ percentage points of the $+24.8568\%$ change in total $B$; built-up land contributes $+7.4416$ points, forest biocapacity $-5.9525$, grazing land $-0.2688$, and fishing grounds $-0.4212$. The non-cropland account sum, which includes built-up land, changes by only $+0.7990$ points ($1.0092\times$) because those unlike account movements net against one another. In the same-edition 1961--2022 sensitivity, the corresponding contributions are cropland $+23.8183$, built-up land $+7.2854$, forest biocapacity $-5.8137$, grazing land $-0.2890$, and fishing grounds $-0.4559$ points; the non-cropland sum changes by $+0.7267$ points ($1.0083\times$). These are accounting attributions, not evidence that forest biomass, soil condition, wetland extent, fish biomass, or modeled $A_c$ declined. Built-up land is an NFA account class, not ecological capital; its inclusion in the non-cropland sum is arithmetic, not ecological aggregation.

\subsection{Physical-stock scoreboard: specification, not inventory}\label{si:stocks}

No common-boundary global physical-stock series for the relevant components was assembled, and no physical-stock composite or global decline finding is presented. This absence does not establish that such a series could never be assembled. Table~\ref{tab:si_stock_scoreboard} is an unpopulated specification of candidate measures and comparability requirements; it contains no empirical endpoint values and is not a model calibration.

Heterogeneous physical stocks cannot be integrated by simply summing their native units: forest biomass or soil carbon (mass), wetland or cropland extent (area), fish biomass by species or assessed stock, freshwater storage (volume), and condition indicators (often dimensionless or categorical) are neither dimensionally commensurate nor necessarily spatially disjoint. Annual flows such as harvest, catch, increment, recharge, or emissions must also remain distinct from stocks and extents. Even after normalization into a common score, overlap, baselines, uncertainty, ecological thresholds, valuation, and permitted substitution determine the interpretation. Component measures therefore form a nonadditive vector unless common spatial and temporal boundaries, normalization or valuation, and substitution rules are declared. A scalar composite is not neutral and cannot be inferred from the NFA gha accounts. This distinction is consistent with ecosystem-extent and ecosystem-condition accounts being reported separately from ecosystem-service and monetary presentations in SEEA Ecosystem Accounting \citep{UN2021SEEA_EA}.

\begin{table}[htbp]
\centering\footnotesize
\caption{Physical-stock scoreboard specification. The rows are unpopulated: no common-boundary global physical-stock series for the relevant components or endpoint values were assembled. Stocks, extents, condition measures, and annual flows are not combined into a composite.}
\label{tab:si_stock_scoreboard}
\begin{tabularx}{\textwidth}{@{}>{\raggedright\arraybackslash}p{0.12\textwidth} >{\raggedright\arraybackslash}p{0.27\textwidth} >{\raggedright\arraybackslash}p{0.22\textwidth} >{\raggedright\arraybackslash}X@{}}
\toprule
Component & Physical stock, extent, or condition to define & Annual flow or pressure (separate) & Comparability rule required \\
\midrule
Forests & Extent (ha); growing stock or biomass (m$^3$ or tonnes C); stated condition & Net annual increment, harvest, or deforestation & Common forest definition and mapped area; reference year; treatment of plantations and overlap \\
Soils & Soil organic-carbon stock or density for a stated depth; structure or erosion condition & Carbon accumulation or loss; erosion where measured & Consistent depth, bulk density, land cover, and sampling method \\
Wetlands & Extent (ha) and an explicit ecological-condition measure & Carbon accumulation, water-storage change, or nutrient flux & Wetland definition and overlap with forest and other land-cover classes \\
Fisheries & Spawning-stock biomass by species or stock (tonnes); age or size structure & Recruitment, catch, and landings & Stock identity, assessment method, reference point, and uncertainty \\
Freshwater & Surface-water or groundwater storage (m$^3$), level, and quality & Recharge and withdrawal & Basin or aquifer boundary; separate stored volume from annual flow \\
\bottomrule
\end{tabularx}
\end{table}

In the two-land model, $A_c$ is an effective area state (ha) and $q$ is a separate capacity-quality index. Neither is directly observed in the NFA aggregate. A physical-stock assessment could inform a model observation operator only after its boundaries and mapping to those effective states are specified. The proxy residual $C=B-X=B-\hat b_fA_f$ is not ecological capital, a measured non-cropland stock, or an estimate of $\Ups(q)A_c$; the NFA components remain account results rather than physical-stock observations.
